\section{Large values from primes above the cutoff}

Fix \(\lambda>1\), and put \(\tau=\lambda w\log w\) and
\(b_p=\exp(\tau/p)-1\). The following product measures the exponential
cost of forcing prime divisors just above \(w\).
\begin{lemma}\label{lem:product}
For \(y=L^2\),
\[
 A_\lambda=\prod_{w<p\le y}\left(1+\frac{b_p}{p}\right),\qquad
 \log A_\lambda=w(\lambda\log\lambda-\lambda+1)+o_\lambda(w).
\]
\end{lemma}
\begin{proof}
Fix \(0<\delta<(\lambda-1)/2\). For \(p=uw\) and
\(1\le u\le\lambda-\delta\),
\[
 \log(1+b_p/p)=(\lambda/u-1)\log w+O_\lambda(1)
\]
uniformly as \(w\to\infty\). The prime number theorem and partial summation
give a contribution
\(w\int_1^{\lambda-\delta}(\lambda/u-1)\dd u+o_{\lambda,\delta}(w)\).
For \((\lambda-\delta)w<p\le(\lambda+\delta)w\), each logarithm is
\(O_\lambda(\log w)\). The prime number theorem bounds their sum by
\(O_\lambda(\delta w)+o_{\lambda,\delta}(w)\). For \(p>(\lambda+\delta)w\),
use \(e^v-1\le ve^v\) for \(v\ge0\). The remaining sum is at most
\[
 \tau w^{\lambda/(\lambda+\delta)}
          \sum_{m>(\lambda+\delta)w}m^{-2}=o_{\lambda,\delta}(w).
\]
Letting \(\delta\) decrease to zero evaluates the integral as stated.
\end{proof}

\begin{lemma}\label{lem:moment}
For either sampling law and each fixed \(\lambda>1\),
\[
 \sum_{i\le K}\E_k e^{\tau T_i(k)}
       \le K A_\lambda e^{o_\lambda(w)},\qquad
 \E_k e^{\tau T_1(k)}\le A_\lambda e^{o_\lambda(w)}.
\]
\end{lemma}
\begin{proof}
When \(p\mid i\), the condition \(p\mid N-iWk\) holds for every \(k\)
if \(p\mid N\), and for no \(k\) otherwise. Every other prime in \((w,y]\)
specifies a single residue class for \(k\bmod p\). In the product
\(\prod_{w<p\le y}(1+b_p\ind_{p\mid n_i})\), take out the constant
factors at \(p\mid(i,N)\) and expand over the primes \(p\nmid i\).
For a squarefree product \(d\), all its congruences specify one class
modulo \(d\) by the Chinese remainder theorem. The sum of all expansion
weights, including the factors forced by \(p\mid(i,N)\), is at most
\[
 \prod_{w<p\le y}(1+b_p)
    =\exp\left(\tau\sum_{w<p\le y}1/p\right)
    \le \exp(O_\lambda(w^2\log w))=Y^{o(1)}.
\]
For integer \(k\), a class contains \(Y/d+O(1)\) points in the interval.
Thus its expected contribution uses \(1/d\), with total additive error
\(Y^{-1+o(1)}\) after normalization by \(Y+O(1)\).

For prime \(k\), take \(d\le\sqrt Y\) first. A non-reduced class
contains no eligible prime, since every prime factor of \(d\) is at most
\(y<Y\). For a reduced class the Brun--Titchmarsh inequality
\cite[Theorem~2, (1.10)]{MV} bounds its probability by
\(O(1/\varphi(d))\), using the prime number theorem for the denominator
and treating \(d=1\) directly. Uniformly over these \(d\),
\(d/\varphi(d)\ll\log y=e^{o(w)}\) by Mertens' theorem. For larger \(d\)
the integer count bounds the probability by
\(O((Y^{-1/2}+Y^{-1})\log Y)\). With the expansion weights, its total
contribution is \(Y^{-1/2+o(1)}\).

Combining the two estimates and retaining all the positive local factors
in \(A_\lambda\), we obtain, under either law,
\[
 \E_k e^{\tau T_i(k)}
   \le e^{o_\lambda(w)} A_\lambda
        \prod_{\substack{w<p\le y\\p\mid(i,N)}}g_p+o(1),
 \qquad g_p=\frac{1+b_p}{1+b_p/p}\le\min(p,1+b_p).
\]
The errors here are uniform in \(i\) and \(N\). To sum the product over
\(i\le K\), expand it in the nonnegative factors \(g_p-1\). A squarefree
product \(d\) divides \(\lfloor K/d\rfloor\) of these indices. Consequently
the average is at most
\[
 \prod_{\substack{w<p\le y\\p\mid N}}\left(1+\frac{g_p-1}{p}\right)
        =e^{o_\lambda(w)} .
\]
Indeed, the factors with \(p\le2\lambda w\) are at most \(2\), and there
are \(O_\lambda(w/\log w)\) of them. The sum of the remaining logarithms
is at most \(\sum_{p>2\lambda w} b_p/p\ll_\lambda \sqrt w\log w\),
by the same exponential bound as in Lemma~\ref{lem:product}. For \(i=1\)
the product over common divisors is empty.
\end{proof}

\begin{proof}[Proof of Theorem~\ref{thm:main}]
Write \(U_i=(\log w)T_i(k)\). For a fixed \(s>0\) take
\(\lambda=e^s\) in Lemmas~\ref{lem:product} and~\ref{lem:moment}.
Markov's inequality yields
\begin{align}
 \E_k\#\{i\le K:U_i\ge s\}
   &\le L^{D+1-e^s+o(1)}, \label{eq:upper}\\
 \Prob(U_1\ge s)&\le L^{1-e^s+o(1)}. \nonumber
\end{align}
Here the exponential penalty is \(\exp(-\lambda sw)\).
Since \(U_i\ge0\), the second inequality gives the limit for the minimum.
For \(s>\log(1+D)\), the first gives the upper bound for the maximum.

For \(1<c<1+D\), let \(P_c=\prod_{w<p\le cw}p\). The prime number theorem
and partial summation give
\[
 \log P_c=(c-1+o(1))w,\qquad
 (\log w)\sum_{w<p\le cw}1/p=\log c+o(1).
\]
The first estimate implies \(P_c=o(K)\). For integer multipliers the
probability of a common factor of \(k\) and \(P_c\) is at most
\(\sum_{w<p\le cw}(1/p+O(1/Y))=O_c(1/\log w)\). Every prime multiplier
in the sampling interval exceeds \(cw\). Whenever \((k,P_c)=1\), the class
\[
             i\equiv N(Wk)^{-1}\pmod{P_c}
\]
contains \(K/P_c+O(1)=L^{D+1-c+o(1)}\) indices in \([1,K]\).
Every prime dividing \(P_c\) divides \(n_i\) at each such index, so
\(U_i\ge\log c+o(1)\). Taking \(c\) arbitrarily close to \(1+D\) from
below gives the lower bound for the maximum.

For a level \(0<t<\log(1+D)\), the same construction with \(c>e^t\)
supplies the lower bound in \eqref{eq:profile} by taking \(c\) as close to
\(e^t\) as the prescribed exponent tolerance requires. Its strict margin
\(\log c-t\) absorbs the error in Lemma~\ref{lem:transfer}. For the upper
bound apply \eqref{eq:upper} at \(s<t\) sufficiently close to \(t\),
then apply Markov's inequality to the count. Lemma~\ref{lem:transfer}
bounds the difference between all the \(U_i\) and \(Z_i\) simultaneously
on a set of probability tending to one. This also transfers the two
extreme-value limits to the representation counts.
\end{proof}
